\documentclass[11pt]{article}
\usepackage[a4paper,margin=1in]{geometry}
\usepackage{fontspec}
\usepackage[boldfont=false]{xeCJK}
\setCJKmainfont{FandolSong-Regular.otf}
\usepackage{amsmath}
\usepackage{amssymb}
\usepackage{array}
\usepackage{booktabs}
\usepackage{graphicx}
\usepackage{adjustbox}
\usepackage{enumitem}
\setlist[itemize]{topsep=2pt,partopsep=0pt,parsep=0pt,itemsep=2pt,leftmargin=1.4em}
\setlist[enumerate]{topsep=2pt,partopsep=0pt,parsep=0pt,itemsep=2pt,leftmargin=1.6em}
\usepackage{parskip}
\usepackage{fvextra}
\fvset{breaklines=true,breakanywhere=true,fontsize=\small}
\setlength{\parindent}{0pt}

\begin{document}

\begin{center}
  {\LARGE\bfseries The macroeconomy}\\[0.35em]
  {\large A Level Economics}
\end{center}
\vspace{0.6em}

\textbf{Macroeconomics} 宏观经济学 looks at the whole economy at once: its total output, its jobs, and its prices. This handout covers how we measure the economy and the main problems it can face.

\section*{National income statistics}

To judge how an economy is doing, we measure its \textbf{national income} 国民收入 — the total value of everything it produces in a year.

\begin{itemize}
\item \textbf{gross domestic product} 国内生产总值 (GDP) is the value of all goods and services made \textbf{inside} a country in a year.

\item \textbf{gross national income} 国民总收入 (GNI) is GDP plus the income a country's people and firms earn \textbf{abroad}, minus income foreigners earn inside it.

\end{itemize}

Two important differences:

\begin{itemize}
\item \textbf{nominal} 名义 values are measured at current prices. \textbf{real} 实际 values are adjusted to remove the effect of rising prices, so they show the true change in output. Always compare real values over time.

\item \textbf{total} values cover the whole country. \textbf{per capita} 人均 values are divided by the population, so they show the average per person.

\end{itemize}

\subsection*{Uses and limitations}

National income is used to measure the \textbf{standard of living} 生活水平 and to compare countries. But it has limits:

\begin{itemize}
\item it ignores the \textbf{distribution} 分配 of income — a high average can hide great inequality.

\item it leaves out unpaid work (housework) and the \textbf{informal economy} 非正规经济 (hidden, untaxed activity).

\item it says nothing about pollution, leisure time, or health.

\end{itemize}

So a higher GDP per capita usually means a better life, but not always.

\section*{The circular flow of income}

The \textbf{circular flow of income} 收入循环流 shows money moving around the economy. \textbf{Households} 家庭 supply factors of production to \textbf{firms} 企业 and receive income (wages, rent, interest, profit). They spend that income on goods and services, and the money flows back to firms.

Money also leaves and enters this flow:

\begin{itemize}
\item \textbf{leakages} 漏出 take money \textbf{out} of the flow: \textbf{saving} 储蓄, \textbf{taxation} 税收, and spending on \textbf{imports} 进口.

\item \textbf{injections} 注入 put money \textbf{into} the flow: \textbf{investment} 投资, \textbf{government spending} 政府支出, and \textbf{exports} 出口.

\end{itemize}

The flow is in \textbf{equilibrium} 均衡 when total injections equal total leakages. If injections are greater, the economy grows; if leakages are greater, it shrinks.

A \textbf{closed economy} 封闭经济 has no trade with other countries. An \textbf{open economy} 开放经济 trades, so it has exports and imports.

\section*{Aggregate demand and aggregate supply}

\subsection*{Aggregate demand}

\textbf{aggregate demand} 总需求 (AD) is the total spending on a country's goods and services at each price level. It has four parts:

\[
AD = C + I + G + (X - M)
\]

\begin{itemize}
\item \textbf{C} is \textbf{consumption} 消费 — spending by households.

\item \textbf{I} is \textbf{investment} — spending by firms on capital.

\item \textbf{G} is \textbf{government spending}.

\item \textbf{(X − M)} is \textbf{net exports} 净出口 — exports minus imports.

\end{itemize}

The AD curve slopes downward: a lower price level means higher real spending, so a greater real output is demanded. AD shifts when any of its four parts changes (for example, a tax cut raises C).

\subsection*{Aggregate supply}

\textbf{aggregate supply} 总供给 (AS) is the total output firms are willing to make at each price level.

\begin{itemize}
\item \textbf{short-run aggregate supply} 短期总供给 (SRAS) slopes upward. In the short run, higher prices make production more profitable, so firms supply more. SRAS shifts when costs change (wages, raw materials, taxes).

\item \textbf{long-run aggregate supply} 长期总供给 (LRAS) shows the economy's full \textbf{productive capacity} 生产能力. It shifts only when the quantity or quality of factors changes.

\end{itemize}

\subsection*{Macroeconomic equilibrium}

\textbf{macroeconomic equilibrium} 宏观经济均衡 is where AD crosses AS. It sets the \textbf{price level} 价格水平 and the \textbf{real output} 实际产出 of the economy. A shift in AD or AS changes both.

\section*{Economic growth}

\textbf{economic growth} 经济增长 is a rise in a country's real output (real GDP) over time.

\begin{itemize}
\item \textbf{actual growth} 实际增长 is the rise in output that really happens (a movement to a point further out, or from inside the LRAS towards it).

\item \textbf{potential growth} 潜在增长 is the rise in the economy's capacity (the LRAS shifting outward).

\end{itemize}

Growth is caused by more or better factors of production: a bigger or more skilled workforce, more investment in capital, new technology, and higher \textbf{productivity} 生产率 (output per worker).

\subsection*{The business cycle}

Real output does not grow smoothly. The \textbf{business cycle} 经济周期 (also called the trade cycle) is the up-and-down pattern of output over the years. It has four phases: \textbf{boom} 繁荣 (fast growth), \textbf{recession} 衰退 (falling output), \textbf{slump} 萧条 (the low point), and \textbf{recovery} 复苏.

\subsection*{Costs and benefits of growth}

\begin{center}
\adjustbox{max width=\linewidth}{%
\begin{tabular}{ll}
\toprule
\textbf{Benefits} & \textbf{Costs} \\
\midrule
higher incomes and living standards & can cause inflation if AD grows too fast \\
more jobs & pollution and use of resources \\
less poverty; more tax for public services & the gap between rich and poor may widen \\
\bottomrule
\end{tabular}}
\end{center}

\section*{Unemployment}

\textbf{unemployment} 失业 means people who are able and willing to work, and are looking for a job, but cannot find one.

\subsection*{Measuring unemployment}

\begin{itemize}
\item the \textbf{claimant count} 申领救济人数 counts only people claiming unemployment benefits. It is cheap but misses those who do not claim.

\item the \textbf{labour force survey} 劳动力调查 asks a large sample whether they are working or seeking work. It follows an international standard, so it is better for comparing countries. The \textbf{labour force} 劳动力 is everyone working or looking for work.

\end{itemize}

\subsection*{Types and causes}

\begin{itemize}
\item \textbf{frictional unemployment} 摩擦性失业 — short-term, while people move between jobs.

\item \textbf{structural unemployment} 结构性失业 — when an industry declines and workers' skills no longer match the jobs available.

\item \textbf{cyclical unemployment} 周期性失业 — caused by low AD in a recession. This is usually the largest type.

\item \textbf{seasonal unemployment} 季节性失业 — when work is only available at certain times of the year (such as farming or tourism).

\end{itemize}

\subsection*{Consequences}

Unemployment wastes output (the economy produces inside its PPC), lowers incomes, and costs the government (more benefits, less tax). It also harms families and can raise crime and ill health.

\section*{Price stability}

\subsection*{Inflation, deflation and disinflation}

\begin{itemize}
\item \textbf{inflation} 通货膨胀 is a sustained rise in the general price level. Money then buys less, so its \textbf{purchasing power} 购买力 falls.

\item \textbf{deflation} 通货紧缩 is a sustained fall in the general price level.

\item \textbf{disinflation} 反通货膨胀 is when inflation is still positive but is slowing down (prices rise, but more slowly).

\end{itemize}

The price level is measured with the \textbf{consumer price index} 消费者价格指数 (CPI). A typical household's spending is used to build a "basket" of goods. Each good is given a \textbf{weight} 权重 for how much is spent on it, and the change in the basket's cost from year to year is the inflation rate.

\subsection*{Causes of inflation}

\begin{itemize}
\item \textbf{demand-pull inflation} 需求拉动型通货膨胀 — AD rises faster than AS, so "too much money chases too few goods".

\item \textbf{cost-push inflation} 成本推动型通货膨胀 — firms' costs rise (wages or raw materials), so they raise prices, shifting SRAS to the left.

\end{itemize}

\subsection*{Consequences}

Inflation lowers purchasing power, especially for people on fixed incomes. It creates uncertainty, can make a country's exports dearer, and helps borrowers but harms savers. Deflation sounds good but can be harmful: people delay buying (waiting for lower prices), so AD falls, firms cut output, and unemployment rises.


\end{document}
